% Vietnamese translation for OI Public Library
% Source-File: IOI中国国家候选队论文/2009/武森/论文ppt/浅谈信息学竞赛中的“0”和“1”.ppt
\documentclass[11pt]{article}
\usepackage{oipl-vn}

\title{Bàn về ``0'' và ``1'' trong thi Tin học\\{\large Ghi chú theo slide}}
\author{Wu Sen\\Trường Trung học số 2 Thạch Gia Trang, Hà Bắc}
\date{2009}

\begin{document}
\maketitle

\origtitle{Zeros and ones in informatics contests}
\sourcepath{IOI-China-Candidate-Team-Papers/2009/Wu-Sen/binary-thinking-slides.ppt}

\section{Chủ đề}

Slide dùng cặp giá trị \(0\) và \(1\) như một lăng kính để nhìn các bài toán:
biểu diễn trạng thái bằng bit, tách một lựa chọn phức tạp thành các quyết định
nhị phân, và khai thác phép toán bit để giảm chi phí.

\section{Ứng dụng trong cấu trúc dữ liệu}

Ví dụ \textit{Matrix} cho thấy cách dùng các bit để biểu diễn trạng thái hàng,
cột hoặc quan hệ bật tắt. Khi trạng thái được nén thành số nguyên, phép xor,
and, or và dịch bit thay thế nhiều vòng lặp nhỏ.

\section{Ứng dụng trong tư duy giải bài}

\begin{itemize}
  \item \textit{Sudoku}: mỗi ô hoặc mỗi tập ứng viên có thể được mô tả bằng
  mặt nạ bit.
  \item \textit{Requirements}: các điều kiện đúng sai được gom thành các biến
  nhị phân, giúp kiểm tra ràng buộc rõ ràng hơn.
  \item \textit{Cow XOR}: trie nhị phân giúp tìm giá trị xor tốt nhất bằng
  cách ưu tiên nhánh đối bit ở từng vị trí.
\end{itemize}

\section{Ghi chú}

Điểm chính của slide là chuyển đổi biểu diễn. Khi bài toán có nhiều lựa chọn
có hoặc không, thuộc hoặc không thuộc, bật hoặc tắt, biểu diễn bit thường làm
cả trạng thái lẫn thao tác trở nên gọn hơn.

\end{document}
